% See https://holatex.app/playground.html# for an online viewer.

\documentclass[tikz]{standalone} % Standalone: Only graphics export.
\usepackage[european,siunitx]{circuitikz}
\usetikzlibrary{arrows.meta}
\usepackage{listings}
\lstset{basicstyle=\footnotesize\ttfamily}
\begin{document}

  \newcommand{\eq}{=}
  \tikzset{ameter/.style={rmeterwa, t=A}}
  \begin{circuitikz}
  %\begin{circuitikz}[american voltages, american currents]

    \tikzset{>=Triangle} % Define tip of arrows drawn with [->]
    \ctikzset{voltage=straight}

    \def\lR{2.5}        % Length of device like R
    \def\dFbY{2.5}        % Delta in direction Y of feedback line
    \def\dRcX{0.3}      % Delta x between RC elements
    \def\dRcY{0.7}      % Delta up/down of the two feedback rows
    \def\dOpX{1.8*\lR}    % Move OP in x to middle of feedback line
    \def\sizeY{\lR}     % Heights Uin
    \def\lUio{0.8}      % Length of I/O connectors
    \def\gapUio{0.3}    % Gap between I/O connector and voltage arrow

    % Input, continued by lower part of feedback line.
    \draw (0,0)
      -- (0,\sizeY) coordinate(in)
      to[R,l=$R$,-*] ++(\lR,0) coordinate(BotLftFb)
      -- ++(0,\dFbY) to[short,-*] ++(\dRcX,0) coordinate(lftRc1)
      -- ++(0,-\dRcY)
      to[C,l=$C_{t_1}$] ++(\lR,0)
      to[short,-*] ++(0,\dRcY) coordinate(rhtRc1)
      to[short,-*] ++(\dRcX,0) coordinate(lftRc2)
      -- ++(0,-\dRcY)
      to[C,l=$C_{t_2}$] ++(\lR,0)
      to[short,-*] ++(0,\dRcY) coordinate(rhtRc2)
      to[short,-*] ++(\dRcX,0) coordinate(lftRc3)
      -- ++(0,-\dRcY)
      to[C,l=$C_{t_3}$] ++(\lR,0)
      to[short,-*] ++(0,\dRcY) coordinate(rhtRc3)
      to[R,l=$R_u$] ++(\lR,0) coordinate(rhtRu);

    % Upper left part of feedback line.
    \draw (lftRc1) -- ++(0,\dRcY)
    to[R,l=$R_{t_1}$] ++(\lR,0) coordinate(rhtRt1)
    -- (rhtRt1 -| rhtRc1) -- (rhtRc1);

    % Upper middle part of feedback line.
    \draw (lftRc2) -- ++(0,\dRcY)
    to[R,l=$R_{t_2}$] ++(\lR,0) coordinate(rhtRt2)
    -- (rhtRt2-|rhtRc2) -- (rhtRc2);

    % Upper right part of feedback line.
    \draw (lftRc3) -- ++(0,\dRcY)
    to[R,l=$R_{t_3}$] ++(\lR,0) coordinate(rhtRt3)
    -- (rhtRt3-|rhtRc3) -- (rhtRc3);

    % OP, output to Ru.
    \draw (BotLftFb) -- ++(\dOpX,0)
      node[plain amp,anchor=-](OP){$OP$}
      (OP.out) -- (OP.out-|rhtRu) coordinate(BotRhtFb)
      -- (rhtRu);

    % Ground line.
    \coordinate(origin) at (0,0);
    \draw (origin) to[short,-*] (origin-|OP.-)
      to[short,-o] ($(origin-|BotRhtFb)+(\lUio,0)$) coordinate(gndOut);

    % OP to ground symbol.
    \draw (OP.+) -- (origin-|OP.-) node[ground]{};

    % Connectors for output voltage.
    \draw (BotRhtFb) to[short, *-o] ++(\lUio,0) coordinate(out);

    % Voltage arrow for output voltage.
    \draw[->] ($(out)+(0,-\gapUio)$) -- node[midway,right]{$U_{OP_q}$}
      ($(gndOut)+(0,\gapUio)$);
  \end{circuitikz}

\end{document}
